\section{Spatial moments without a positive transfer spectrum}\label{sec:transfer}
Fix $\beta>0$, $g\ge0$, and an even transverse width $m\ge4$. Group the real representative \eqref{eq:realH}
into columns. The column Hilbert space is finite; let $h_m$ contain its rung
and vertical interactions, and let $G_b$ be the real insertions on its sites.
Then the strip Hamiltonian has the form
$\sum_j h_{m,j}-g\sum_{j,b}G_{b,j}G_{b,j+1}$.
The width dependence of $h_m$, the insertions, and the transfer dimension is
allowed. No bounded-width approximation is taken.

For $q\ge1$ time slices put $\epsilon=\beta/q$, and let an event label be either
$\varnothing$ or one insertion label $b$. Define the symmetric two-history
vertex matrix
\begin{equation}\label{eq:vertex}
 W_{\epsilon}(x,y)=
 \begin{cases}
 I-\epsilon h_m,&x=y=\varnothing,\\
 \sqrt{\epsilon g}\,G_b,&\{x,y\}=\{\varnothing,b\},\\
 0,&x\ne\varnothing,\ y\ne\varnothing.
 \end{cases}
\end{equation}
For histories $\xi,\eta\in(\{\varnothing\}\cup\mathcal B)^q$, define
\begin{equation}\label{eq:transfer}
 (X_q)_{\xi\eta}=
 \sum_{a_0,\ldots,a_{q-1}}\prod_{t=0}^{q-1}
   W_{\epsilon}(\xi_t,\eta_t)_{a_t,a_{t+1}},\qquad a_q=a_0.
\end{equation}

\begin{proposition}[Constructed finite spatial moment limits]\label{prop:transfer}
Each $X_q$ is real symmetric. For each fixed even physical circumference $n\ge4$,
\begin{equation}\label{eq:momentlimit}
 \lim_{q\to\infty}\tr X_q^n=Z_{n,m}(\beta).
\end{equation}
The same statement holds with the two spatial directions interchanged.
\end{proposition}
\begin{proof}
The symmetry of \eqref{eq:vertex} in its two event labels proves
$(X_q)_{\xi\eta}=(X_q)_{\eta\xi}$ directly. Neither entrywise positivity nor
transposition of $G_b$ is used. Contracting the same finite rectangular tensor
network first in time or first in space gives
$\tr X_q^n=\tr P_n(\epsilon)^q$. Here $P_n$ is a finite matching polynomial:
occupied bonds at a time slice cannot share a vertex, and unoccupied vertices
supply $I-\epsilon h_m$. Its constant term is $I$ and its first derivative is
$-H_{n,m}$. Every selected edge has two square-root factors whose product is
$\epsilon g$. Thus $P_n(t)=I-tH_{n,m}+O(t^2)$ in a finite matrix algebra.
A noncommutative telescoping estimate for powers implies
$P_n(\beta/q)^q\to e^{-\beta H_{n,m}}$ in norm. Trace continuity proves
\eqref{eq:momentlimit}. The column-to-torus reindexing is an exact equality of
onsite and edge lifts. Swapping the two spatial axes gives the second direction.
\end{proof}

The formal route uses these finite symmetric approximants and their trace
limits. It does not need to construct a single infinite-dimensional limiting
transfer operator or prove its spectral completeness.
For even $n$, if $\lambda_{q,i}$ are the real eigenvalues of $X_q$ and the
denominator is nonzero, the weights
\begin{equation}\label{eq:spatialweights}
 p_{q,i}=\frac{\lambda_{q,i}^{n}}{\sum_j\lambda_{q,j}^{n}}
\end{equation}
are nonnegative and sum to one, even if some $\lambda_{q,i}<0$.
The denominators are positive eventually, since their limits are the strictly
positive physical partition functions. Consequently
$\tr X_q^{2n}/(\tr X_q^n)^2$ is a genuine probability purity.
Passing the finite inequalities to the limit legitimizes the spatial ratios
in \eqref{eq:staircase}. Thermal purity is already the purity of the finite
Gibbs weights $e^{-\beta E_j}/Z(\beta)$.

This argument isolates the two-dimensional issue: each fixed transverse width
has its own transfer spectrum, but the probability inequalities used below
have constants independent of its dimension or width. No common eigenbasis,
joint $x/y$ spectrum, or algebraic relation between different widths is assumed.
